\documentclass[10pt,a4paper]{amsart}
\usepackage[margin=19mm]{geometry}
\usepackage{amsmath,amssymb,mathtools}
\usepackage[scr=rsfs,cal=euler]{mathalfa}
\usepackage{xfrac}
\usepackage[unicode,hidelinks,bookmarks=false]{hyperref}
\newcommand{\ie}{i.e.\ }
\newcommand{\cf}{cf.\ }
\newcommand{\resp}{respectively\ }
\newcommand{\Subsection}{\subsection}
\providecommand{\nobreakdash}{\nobreak\hbox{-}}
\setlength{\emergencystretch}{2em}
\multlinegap0pt
\usepackage{bm}
\usepackage{bbm}
\usepackage{dsfont}
\usepackage{xspace}
\usepackage{cancel}
\usepackage{mathtools}
\usepackage{amsthm}
\makeatletter
\@ifundefined{theorem}{\newtheorem{theorem}{Theorem}}{}
\@ifundefined{defn}{\newtheorem{defn}[theorem]{Definition}}{}
\@ifundefined{prop}{\newtheorem{prop}[theorem]{Proposition}}{}
\makeatother
\DeclareMathOperator*{\argsort}{arg\,sort}
\newcommand{\RR}{\mathbb{R}}
\newcommand{\abs}[1]{\left\lvert#1\right\rvert}
\title[Bi-Lipschitz Ansatz for Anti-Symmetric Functions]{Bi-Lipschitz Ansatz for Anti-Symmetric Functions}
\author{Nadav Dym, Jianfeng Lu, Matan Mizrachi}
\date{Source: arXiv:2503.04263v2 (CC BY 4.0)}
\begin{document}
\maketitle

\noindent\emph{Source and licence.} Bi-Lipschitz Ansatz for Anti-Symmetric Functions. Version arXiv:2503.04263v2 (CC BY 4.0), distributed under the Creative Commons Attribution 4.0 International licence, as recorded in the arXiv licence metadata for this exact version and shown on the abstract page https://arxiv.org/abs/2503.04263v2. Full rights notice: licenses/arxiv-2503.04263-CC-BY-4.0.txt.\par\smallskip

\noindent\textit{Editorial scope.} Counted unit: the Prop whose source label is \texttt{prop:stab\_operator\_correct} in the article (source0.tex lines 704--706, source label \texttt{prop:stab\_operator\_correct}), delivered together with the article's own complete original proof of it (source0.tex lines 707--772, one \texttt{proof} environment of 66 lines). No mathematics is written, repaired or re-derived: statement and proof are the article's own text line for line. Required context retained uncounted: def:afa (lines 551--576); eq:tau (lines 661--662); eq:w (lines 664--666); def:staboperator (lines 689--692). Every definition, display and label of those lines is retained unchanged. Omitted from this package and not redistributed: 1--550, 577--660, 663--663, 667--688, 693--703, 773--1226. Nothing inside a retained excerpt is omitted or altered: every retained body is delivered exactly as the article's lines are written, so no omission receipt of any kind accompanies this package. Citations: the retained excerpts carry no citation command at all, so no bibliography entry of the article is transcribed and no reference is redistributed. Reference adaptation: none was needed and none was made; every reference inside the delivered unit resolves inside it, so no unresolved reference or citation is produced. Printed slips kept literal: no printed word, symbol, hypothesis or number of the counted statement or of its proof was altered. Layout and class: the article's own class is \texttt{amsart} with packages including amsmath, amssymb, amsthm, amsrefs, indentfirst, mathrsfs, parskip, graphicx, wrapfig, hyperref, xcolor, fourier, geometry, booktabs; the wrapper is the lane's pinned \texttt{amsart} editorial wrapper at 10pt on A4, which restates the theorem-like environments the retained text needs and adds the article's own macro definitions that the retained text needs (\texttt{\textbackslash RR}, \texttt{\textbackslash abs}, \texttt{\textbackslash argsort}), with extra wrapper packages (bm, bbm, dsfont, xspace, cancel, mathtools). The delivered numbering is the unit's own. Assets: no retained span contains a figure environment, a graphics-inclusion command, a PDF-page inclusion, a table or a TikZ picture; no image file is copied into this package, the package declares no asset dependency and no image file is redistributed with it. Licence: version arXiv:2503.04263v2 of this article is distributed under the Creative Commons Attribution 4.0 International licence, as recorded in the arXiv licence metadata for this exact version and shown on the abstract page https://arxiv.org/abs/2503.04263v2. A full rights notice is delivered at licenses/arxiv-2503.04263-CC-BY-4.0.txt. Characters: every retained excerpt is pure ASCII, so the delivered engine, XeLaTeX with its default Latin Modern text font, reports no missing character.
\begin{defn}[WS and Antisymmetric Frame Averaging Operator] \label{def:afa}
Let $d,n,m$ be natural numbers, and let $\mathcal{F}_{\text{as}}$ be a weighted permutation frame of size $m$. Let  $\Phi[\cdot]$ be a linear operator mapping functions on $\RR^{d\times n}\setminus \Omega_{d\times n}$ to functions on the same space. We say that $\Phi $ is a \textit{weakly stabilizing (WS) operator} for $\mathcal{F}_{\text{as}}$, if
\begin{enumerate}

\item For every compact $K \subset \RR^{d \times n}$, which is $S_n$-invariant, there exists a positive constant $C_K$, such that
  \begin{equation} \label{eq:cond_bound_antis}
    \forall f\in  C\left(\RR^{d \times n}\right), \quad   \max_{x \in K}{\abs{\Phi[f](x)}} \leq C_K \max_{x \in K}{\abs{f(x)}}
  \end{equation}
  \item If $f:\RR^{d\times n} \to \RR$ is antisymmetric, then $\Phi[f] = f$.
    \item \label{cond:proj} For every $f\in  C\left(\RR^{d \times n} \setminus \Omega_{d\times n} \right)$, the function $\Phi[f]$ will also be continuous on $\RR^{d\times n}\setminus \Omega_{d\times n} $.
    \item For any convergent sequence 
  $x^k \to x$ in $\RR^{d \times n} \setminus{\Omega_{d \times n}}$, with limit $x \in \Omega_{d \times n}$,
  \begin{equation} \label{eq:cond_weak_convergence}
\forall f\in  C\left(\RR^{d \times n}\right), \quad     \Phi[f]\left(x^k \right)  \stackrel{k \to \infty}{\rightarrow} 0
  \end{equation}

\end{enumerate}
An \emph{antisymmetric frame averaging operator} will then be an operator
\begin{equation} \label{eq:form2frame}
    \begin{split}
        \mathcal{E}[f](x) &\coloneq 0, \quad \forall x \in \Omega_{d \times n}, \\
        \mathcal{E}[f](x) &\coloneq \sum_{i=1}^{m}{w_i(x)(-1)^{\tau_i}\Phi[f]\left( \tau_i^{-1} x \right)}, \quad \forall x \in \RR^{d \times n} \setminus{\Omega_{d \times n}}.
    \end{split}
    \end{equation}
    defined by a weighted permutation frame $\mathcal{F}_{\text{as}}$ and an appropriate WS operator $\Phi$. 
\end{defn}

\begin{equation} \label{eq:tau}
	 \tau_i^{-1}(x) \coloneq \argsort{\left(x_1 \cdot a_i, \dots, x_n \cdot a_i \right)}  .\end{equation}

\begin{equation}\label{eq:w}
	w_i(x) \coloneq \frac{\abs{Q\left(x \cdot a_i \right)}}{\sum_{l \in [m]}{\abs{Q\left(x \cdot a_l \right)}}}=\frac{\min_{1\leq j<k\leq n} |a_i \cdot(x_j-x_k) |  }{\sum_{l\in[m]}\min_{1\leq j<k\leq n} |a_l \cdot(x_j-x_k) | },
\end{equation}

\begin{align} \label{def:staboperator}
        \Phi[f](x) &\coloneq f(x) - \sum_{\sigma_{i,j} \in H}{\omega_{i,j}(x)\left(f(x) + f(\sigma_{i,j} x) \right)}, \quad \forall x\in \RR^{d\times n}\setminus \Omega_{d\times n}\\
         \Phi[f](x) &\coloneq 0, \quad \forall x \in \Omega_{d\times n} \nonumber
\end{align}

\begin{prop} \label{prop:stab_operator_correct}
Given the weighted permutation frame $\mathcal{F}_{\text{as}} $ defined at \eqref{eq:tau} and \eqref{eq:w}, the operator $\Phi[\cdot]$ defined in \eqref{def:staboperator} is a  \textit{WS} operator.
\end{prop}

\begin{proof}
	 We will show that all related conditions in the definition \ref{def:afa} hold for $\Phi$:
	 
	\textbf{$\Phi$ is Bounded} 
For every compact $K \subset \RR^{d \times n}$, which is $S_n$-invariant, and for any function $f: \RR^{d \times n} \to \RR$, we bound $\Phi[f]$ over $K$, by noting the for all $x \in  \Omega_{d\times n}$ we have $\Phi[f](x)=0 $, and for all $x \in \RR^{d\times n}\setminus \Omega_{d\times n}$,	
\begin{equation}
		\begin{split}
			\max_{x \in K}{\abs{\Phi[f](x)}} 
            &= 
			\max_{x \in K}{\abs{f(x) - \sum_{\sigma_{i,j} \in H}{\omega_{i,j}(x)\left(f(x) + f(\sigma_{i,j} x) \right)}}}
			\\
			&=
			\max_{x \in K}{\abs{\left(1 - \sum_{\sigma_{i,j} \in H}{\omega_{i, j}(x)}\right)f(x) - \sum_{\sigma_{i,j} \in H}{\omega_{i,j}(x) f(\sigma_{i,j} x)}}}
			\\
			&\leq 
			\frac{1}{2}\max_{x \in K}{\abs{f(x)}} + \sum_{\sigma_{i,j} \in H}{\omega_{i,j}(x) \max_{x \in K}{\abs{f(\sigma_{i,j} x)}}}
			\\
			&\leq \frac{1}{2}\max_{x \in K}{\abs{f(x)}} + \frac{1}{2}\max_{x \in K}{\abs{f(x)}} = \max_{x \in K}{\abs{f(x)}}
		\end{split}
	\end{equation} 
	\textbf{Antisymmetry}
	  Let $f:\RR^{d\times n}\to \RR$ be an antisymmetric function. Then, for all $x \in \RR^{d \times n}$, and for all $\sigma_{i,j} \in H$, we have $f(x) = -f(\sigma_{i,j} x)$, and so $\Phi[f] = f$. 
	
	\textbf{Continuity in $\RR^{d\times n}\setminus \Omega_{d\times n}$} The function $\eta: \RR^{d \times n} \to \RR$ is continuous everywhere, and so is the cotangent function over the domain $(0, \frac{\pi}{2}]$. Accordingly, for all $x \in \RR^{d \times n} \setminus{\Omega_{d \times n}}$, and for any fixed $\delta > 0$, we have $\frac{\pi\min{\left\{\left\| x_i - x_j \right\|_2, \delta + \eta(x) \right\}}}{2\eta(x) + 2\delta} \in (0, \frac{\pi}{2}]$ ensuring the continuity of $\tilde{\omega}_{i, j}$ for all $1 \leq i < j \leq n$. And so, for every $f\in  C\left(\RR^{d \times n} \setminus \Omega_{d\times n} \right)$, the function $\Phi[f]$ will also be continuous on $\RR^{d\times n}\setminus \Omega_{d\times n} $, as summing over $\sigma_{i, j} \in H$ the multiplication of the function $f(x)$, the functions $f(\sigma_{i,j}x)$ and the weight functions $\omega_{i,j}(x)$, which are all continuous outside of $\Omega_{d \times n}$, is continuous.
	
	\textbf{Continuous limit}
    Let $x^k \to x$ be any convergent sequence in $\RR^{d \times n} \setminus{\Omega_{d \times n}}$, with limit a $x \in \Omega_{d \times n}$ at which $\eta(x) = 0$. 
    To simplify, we define the complementary subsets $I^+ ,I^- \subseteq H$ by
    \begin{equation}
    \begin{split}
        I^+ &\coloneq \left\{i, j \in [n]: x_i = x_j, i \ne j \right\} \ne \emptyset, \\
        I^- &\coloneq H \setminus{I^+}.
    \end{split}
    \end{equation}
For all $(i,j)\in I^+$, we will have that $\|x_i^k-x_j^k\| \rightarrow 0 $, which will imply that $\tilde\omega_{i,j}(x^k)\rightarrow \infty $. In contrast, for every $(i, j) \in I^-$, the weight $ \tilde\omega_{i,j}(x^k) $ has a (finite) limit, and therefore, for every $(i, j) \in I^-$
\begin{equation}\label{eq:small_omega}
\begin{split}
    &\lim_{k \to \infty}{\omega_{i,j}(x^k)} \\  &= 
 \frac{1}{2} \lim_{k \to \infty}{
\frac{\tilde{\omega}_{i,j}(x^k)}{\sum_{(l, k) \in I^+}{\tilde{\omega}_{l,k}(x^k)}  + \sum_{(l, k) \in I^-}{\tilde{\omega}_{l,k}(x^k)}
} 
} \\ &\leq
 \frac{1}{2} \lim_{k \to \infty}{
\frac{\tilde{\omega}_{i,j}(x^k)}{\sum_{(l, k) \in I^+}{\tilde{\omega}_{l,k}(x^k)}
} 
} = 0.
\end{split}
\end{equation}
Now, let $f \in C(\RR^{d \times n})$ be any continuous function. By definition, the sequence $\abs{f(x^k) - f(\sigma_{i, j} x^k)}$ converges to zero for all $(i, j) \in I^+$ for any such function $f$. Finally, considering \eqref{eq:small_omega}, we hereby conclude the continuity of $\Phi[\cdot]$ with the limit
\begin{equation}
\begin{split}
     &\lim_{k \to \infty}{\left(
f(x^k) - \sum_{\sigma_{i,j} \in H}{\omega_{i,j}(x^k)\left(f(x^k) + f(\sigma_{i, j} x^k) \right)}
\right)} \\ &= 
\lim_{k \to \infty}{\left(
f(x^k) - \sum_{\sigma_{i,j} \in I^+}{\omega_{i,j}(x^k)\left(f(x^k) + f(\sigma_{i, j} x^k) \right)}
\right)}
\\ &= 
\lim_{k \to \infty}{\left(
f(x^k) - 2\sum_{\sigma_{i,j} \in I^+}{\omega_{i,j}(x^k)f(x^k)}
\right)} \\
&= \lim_{k \to \infty}{\left(f(x^k) - f(x^k) \right)} = 0
\end{split}
\end{equation}
as desired, to show that $\Phi[\cdot]$ is a linear WS operator.
\end{proof}

\end{document}
